% answers_short.tex —— 2008 年·普通高等学校招生考试·陕西卷·文科 · 简短答案速查
% 由 tools/make_short_answers.py 自 answers.tex 的 \gksolhead/\gkanswer 提取，
% 并与 2008/2008_数学真题_陕西文_答案要点.md 逐题比对后生成，勿手改（改源 MD 后重跑）。

\begin{gknotice}
\textbf{使用说明}：本节只给最终答案便于快速核对，完整推导见解析版。答案为编写者独立推导，\textbf{非官方标准答案}。
\end{gknotice}

\gkpart{选择题}
\begin{center}\begin{tabular}{ccccccccccccc}
\toprule
题号 & 1 & 2 & 3 & 4 & 5 & 6 & 7 & 8 & 9 & 10 & 11 & 12 \\
\midrule
答案 & B & D & C & B & A & A & D & C & B & D & A & C \\
\bottomrule
\end{tabular}\end{center}

\gkpart{填空题}
\begin{itemize}[leftmargin=2.6em]
  \item[\textbf{13.}] $\sqrt2$。
  \item[\textbf{14.}] $84$。
  \item[\textbf{15.}] ②。
  \item[\textbf{16.}] $96$。
\end{itemize}

\gkpart{解答题}
\begin{itemize}[leftmargin=2.6em]
  \item[\textbf{17.}] （1）$T=4\pi$，最大值 $2$，最小值 $-2$；（2）偶函数。
  \item[\textbf{18.}] （1）$\dfrac16$；（2）$\dfrac7{12}$。
  \item[\textbf{19.}] （1）证明见解析；（2）$\arccos\dfrac{\sqrt{21}}7$。
  \item[\textbf{20.}] （1）证明见解析；（2）$S_n=\dfrac{n(n+1)}2+2-\dfrac{n+2}{2^n}$。
  \item[\textbf{21.}] （1）证明见解析；（2）$k=\pm2$。
  \item[\textbf{22.}] （1）增 $(-\infty,-a)$、$\left(\dfrac a3,+\infty\right)$，减 $\left(-a,\dfrac a3\right)$；（2）$\left(-\infty,1-\dfrac{\sqrt2}2\right]$；（3）$a\leqslant-3$ 或 $a>1$。
\end{itemize}
